% Source: https://artofproblemsolving.com/community/c3h2520453_algebra_b_claims
% Stable identity: aops-topic-2520453-q01
For $x\ne0$ factor out the highest powers:
\[p(x)=x^m\bigl(a+u(x)\bigr),\qquad q(x)=x^n\bigl(b+w(x)\bigr),\]
where $u(x)$ is a sum of terms $\dfrac{c}{x^j}$ with $j\ge1$ (coming from the lower-degree terms of $p$), and similarly for $w(x)$. Each such term tends to $0$ as $x\to\pm\infty$, so $u(x)\to0$ and $w(x)\to0$. Hence, for large $|x|$,
\[r(x)=x^{m-n}\cdot\dfrac{a+u(x)}{b+w(x)},\qquad\text{where}\qquad\dfrac{a+u(x)}{b+w(x)}\to\dfrac ab\ne0.\]
\textbf{(a)} If $m<n$, then $x^{m-n}=\dfrac1{x^{n-m}}\to0$, and a quantity tending to $0$ times one tending to $\dfrac ab$ tends to $0$.

\textbf{(b)} If $m=n$, then $x^{m-n}=1$ and $r(x)\to\dfrac ab$, the ratio of the \emph{leading} coefficients.

\textbf{(c)} If $m>n$, then $\left|x^{m-n}\right|\to\infty$ while the second factor stays close to $\dfrac ab\ne0$ (for large $|x|$ its absolute value exceeds $\dfrac12\left|\dfrac ab\right|$). So $|r(x)|\to\infty$, and $r(x)$ cannot approach any number $L$.

\textbf{(d)} The function is \emph{defined} by the formula $\dfrac{x^2-9}{x-3}$, which at $x=3$ asks us to divide $0$ by $0$; so $3$ is not in the domain. The cancellation $\dfrac{(x-3)(x+3)}{x-3}=x+3$ is only valid when $x-3\ne0$. Thus the graph of $f$ is the line $y=x+3$ with the single point $(3,6)$ removed: a hole. (Near $x=3$ the values approach $6$, but $f(3)$ itself does not exist.)
